\section{Related Work}
\label{sec:related}

\paragraph{LLM-as-judge bias and construct validity.} A growing body of
work questions whether LLM judges measure what evaluators intend
\citep{gu2024survey}. \citet{chen2024humans} directly document gender,
authority, and beauty bias in both human and LLM judges and show these
biases are exploitable via targeted attacks, but treat gender as one bias
type among several general-purpose evaluation tasks rather than isolating
it with a factorial design in a single, fixed domain. A related but
distinct use of LLM judges appears in \citet{kumar2024decoding}, who use
an LLM as the instrument to \emph{detect} gender bias in the outputs of
other models; there, bias lives in the judged model's generations, whereas
the present study asks whether the bias lives in the judge itself.

\paragraph{AITA and moral-judgment benchmarks.} Scruples
\citep{lourie2021scruples} established the WHO-task reframing of AITA-style
disputes adopted here (Section~\ref{sec:dataset}). More recently,
\citet{sachdeva2025normative} evaluate seven LLMs' blame judgments
against human verdicts on over 10,000 real AITA dilemmas, finding each
model is internally self-consistent but that models diverge substantially
from one another and from human judgments -- motivating this paper's RQ2
(cross-model heterogeneity) and its choice to stratify by model rather
than pool across models. Neither of these prior AITA-derived evaluations,
however, manipulates the disputants' gender under a matched, factorial
design; the underlying vignettes are naturally occurring and confounded
with whatever demographic signal happened to be present in the original
post.

\paragraph{Gender and relational moral judgment in humans.}
\citet{betti2025moral} find, using a large corpus of online moral
discourse, no general causal effect of a story's protagonist's gender on
the moral judgment they receive from human commenters -- \emph{except} in
relationship and friendship stories specifically, where male protagonists
received more negative evaluations than female protagonists in matched
comparisons. This is a direct human-judgment precedent for exactly the
domain this paper studies (relationship-norm disputes), and motivates
asking whether an analogous effect transfers to LLM judges performing the
same kind of evaluation, or is instead a distinctly human phenomenon.

\paragraph{Closest prior work: gender asymmetry in LLM moral judgment.}
Two prior studies directly anticipate RQ1, in opposite directions.
\citet{bajaj2024genmo} introduce GenMO, parallel short stories differing
only in a character's gender, and find several production models
(including Claude 3 Sonnet and Opus) show a consistent gender-swap effect,
inclining toward the female character in a majority of cases -- a
general moral-judgment result, not specific to relationship-conflict
framing, and not broken down by norm-violation type. More directly
relevant, \citet{si2026gama} introduce GAMA-Bench, a gender-mirrored
benchmark of intimate-relationship and public-conflict scenarios scored
on six response-framing dimensions (punitive wording, therapeutic
wording, severity, empathy toward the actor, accusatory content, and
full-blame attribution) across ten current models. On the intimate track
specifically, they find a consistent \emph{male-disadvantaging} asymmetry
-- male actors receive more punitive, blame-centered framing for matched
misconduct -- that persists across model families, model scale, and
reasoning mode, and is not reduced (and sometimes widened) by
chain-of-thought reasoning. \citet{si2026gama} is the closest existing
work to this paper's design in spirit (matched gender-mirrored
relationship-conflict scenarios, structured multi-dimensional scoring
rather than a single verdict), and its finding must be engaged with
directly rather than treated as background: it already establishes a
relationship-conflict-specific gender asymmetry, in the opposite direction
from \citet{bajaj2024genmo}'s general-domain result, that persists across
model scale and reasoning style. This paper's distinguishing contribution
is the decomposition by discrete relationship-norm-violation family (9
families: household labor, childcare, financial provision, jealousy, and
so on) using a single fixed fault-rating scale rather than open-ended
response-framing metrics -- RQ2 asks directly whether the direction and
magnitude of any gender asymmetry \emph{varies by which norm is violated},
a question none of these three studies (Betti et al., Bajaj et al., Si
et al.) answers, since none decomposes by norm-violation type.

\paragraph{Hedging and both-sides validation.} \citet{cheng2025elephant}
document ``social sycophancy'' in LLMs, including a tendency to affirm
both parties in a paired moral-conflict scenario as not at fault in
roughly half of cases -- direct evidence for the concern motivating RQ3
and the hedge-rate tracking described in Section~\ref{sec:protocol}: a
model's willingness to render a judgment at all, and whether that
willingness itself varies by agent gender, is a site where bias could
hide from any analysis that only looks at committed ratings.

% TODO: this section merges two independent literature passes (2026-08-11):
% an initial pass and the user's own separately-compiled lit_review.md
% (paper/sources/), cross-checked against that doc's own verification
% audit before any citation was added here. Still not checked against the
% workshop's own accepted-papers/related submissions, and should be
% revisited once Section~\ref{sec:results} exists in case specific
% findings there change which prior work is most relevant to foreground.
% Background-only sources from lit_review.md (care-ethics grounding,
% marital-attribution literature beyond the one construct-precedent
% citation in measurement_protocol.tex, general contamination-detection
% vocabulary) were deliberately left out of this section given the 4-page
% short-paper budget -- paper/sources/lit_review.md remains the fuller
% reference if space allows expansion later.
